\section{Introduction}\label{introduction}

Food inflation is not only a question of the expected change in a price
index. Households, governments, humanitarian agencies, and firms are
also exposed to the upper tail of future price changes: the possibility
that food prices rise much faster than their central forecast. This
distinction is especially important in economies where household
expenditure is concentrated on food and where exchange-rate, fuel-price,
transport, security, and supply shocks can cause rapid changes in retail
prices. A distributional forecast can therefore be more
decision-relevant than a point forecast, because it directly measures
the severity of plausible adverse outcomes.

The inflation-at-risk literature formalises this idea by modelling
conditional quantiles of future inflation rather than only its
conditional mean (\citeproc{ref-lopezsalido2024inflation}{López-Salido
and Loria 2024}; \citeproc{ref-banerjee2024inflation}{Banerjee et al.
2024}). Existing applications largely concern aggregate consumer-price
inflation in advanced and emerging economies. Food prices create a
harder forecasting problem. They are highly disaggregated, seasonal,
heterogeneous across commodities and locations, and vulnerable to
measurement noise. More importantly, the upper tail may shift quickly
after a macroeconomic or supply-chain regime change. A model estimated
in a stable historical period can be well calibrated in ordinary months
and still fail severely when a common upward shift moves most
state-commodity observations beyond the training distribution.

This study examines that failure mode in Nigeria. The empirical setting
contains a major distribution shift during 2023, when exchange-rate and
fuel-price conditions changed sharply and food-price inflation became
persistently elevated. The forecasting problem is operationally
constrained: the target is three-month forward inflation, so the error
attached to a forecast issued in month t is not fully observable until
t+3. An adaptive model must therefore respond to an ongoing break using
current covariates and stale forecast errors. This delayed-feedback
structure differentiates the problem from conventional rolling-window
evaluation.

The paper makes four contributions. First, it constructs a granular Food
Inflation-at-Risk target from a panel of 37 Nigerian states/FCT and 43
commodities and documents the collapse of static upper-tail calibration
during the 2023 regime change. Second, it evaluates several transparent
delayed-feedback adaptations: exponentially weighted common-shift
adjustment, rolling conformal calibration, and hierarchical partial
pooling across commodities and markets. Third, it develops an inertial
regime-aware controller that combines these expert forecasts using
observable stress diagnostics and only forecast losses that have matured
after the three-month target horizon. Fourth, it evaluates temporal
persistence, cross-source transport, subgroup calibration, and simulated
structural breaks rather than relying on one favourable test period.

The headline result is deliberately restrained. The meta-controller has
lower pooled post-2022 pinball loss than every individual expert in both
the same-source NBS continuation and the external WFP market panel. It
reduces loss significantly relative to a static quantile, but its
advantage over the best expert is not statistically decisive. It also
inherits a recovery problem: the stress override remains too
conservative during the 2025 NBS extension. The evidence therefore
supports regime-aware forecast combination under delayed feedback, not a
claim that the proposed controller is universally superior.

The remainder of the paper is organised as follows. Section 2 positions
the study within the inflation-at-risk, conformal prediction, and online
expert-aggregation literatures. Section 3 describes the data, outcome,
chronological splits, and information constraints. Section 4 defines the
expert forecasts and meta-controller. Section 5 reports the
development-panel failure, extended and external validation, regret
analysis, subgroup coverage, and simulations. Section 6 discusses
interpretation and limitations. Section 7 concludes and states the
frozen confirmatory-validation requirement.
